\section{Threats to Validity}
\subsection{Construct validity}
``Attribution correctness'' depends on the ground-truth definition. A semantic tool can indirectly cause a process through wrappers or schedulers, making immediate parentage different from causal responsibility. The evaluation must therefore define relation types precisely---e.g., direct execution, delegated execution, shared-resource use---rather than reduce every relation to one binary edge. Ground-truth manifests should be independently generated and audited.

Feature comparisons have another construct threat: a general tracing platform can often represent custom process/file/network spans even if it does not collect them natively. We therefore use ``not native'' rather than ``unsupported'' when official documentation lacks an independent OS collector, and controlled experiments should distinguish default instrumentation from deliberately extended custom instrumentation.

\subsection{Internal validity}
Provider integrations evolve quickly. A result measured with one release may reflect an integration bug rather than a conceptual limitation. All comparative experiments must freeze versions, configuration, model endpoints, and workload seeds/fixtures. For nondeterministic model-driven workloads, the OS side effects should be induced by deterministic tool scripts after the agent chooses a known action, or the study should replay a fixed semantic trace into the execution harness.

Timing-sensitive runtime tests can also be flaky. Blind-window experiments should model observation schedules deterministically where possible and complement real-time stress tests rather than rely only on ``sleep for $n$ milliseconds'' assumptions.

\subsection{External validity}
The supported coding-agent and local-runtime integrations do not represent every agent architecture. Browser agents, mobile agents, Kubernetes-hosted agents, serverless tool execution, and remote sandboxes introduce different visibility boundaries. The proposed evidence model should generalize, but the implementation and measurements may not. Similarly, developer workstations differ from production clusters in workload scale and privilege constraints.

\subsection{Conclusion validity}
Performance distributions can be heavy-tailed, especially when model/network latency dominates. Relative overhead should therefore be measured on deterministic local workloads as well as realistic agent sessions, report medians and tail statistics, and include sufficient repetitions. Human-study claims require power analysis and appropriate mixed-effects or repeated-measures statistics rather than relying on a small convenience sample.

\section{Reproducibility, Ethics, and Data Handling}
The public implementation should be accompanied by a frozen artifact containing workload generators, environment manifests, competitor configuration instructions that comply with each system's license and service terms, raw metric outputs, and scripts that compute reported statistics. The manuscript itself should contain static final numbers and must not depend on executing external data-loading code during typesetting.

Observability studies can capture sensitive prompts, source code, authentication material, file contents, and network destinations. Experimental workloads should therefore use synthetic repositories and controlled endpoints whenever possible. Human-subject evaluation should avoid collecting participants' real agent histories. Any released trace corpus should be generated for the study or reviewed for secrets before publication.

A further ethical issue is attribution. A graph that names an agent or process can be interpreted as assigning responsibility. The interface should preserve evidence grade, inference status, and ambiguity so that derived correlation is not presented as forensic certainty. This is both a systems-design requirement and a reporting requirement for the research artifact.

\section{Future Research Directions}
The present design opens several research questions beyond the initial system paper. First, propagated cryptographic or trace identities could reduce ambiguity without relying on command matching, but require adoption across agent, tool, gateway, and model-runtime boundaries. Second, learned correlation could potentially improve recall, yet any learned model would need calibrated abstention and auditable evidence features to avoid replacing a transparent heuristic with opaque false attribution. Third, whole-system provenance backends such as LSM/eBPF-style collectors could strengthen Linux evidence while retaining the canonical cross-layer model. Fourth, distributed agents require provenance across machines and trust domains, bringing the problem closer to X-Trace and distributed provenance but with agent identities as first-class principals. Finally, graph-layout evaluation could become a separate human-factors study focused on shared-resource hubs, semantic group continuity, and dynamic orientation rather than generic crossing counts.

\section{Conclusion}
Agent observability is rapidly becoming graph-based, but a graph of application spans is not the same object as a graph that connects semantic agent actions to independently observed machine effects. \system treats this distinction as the central problem. It preserves provider semantics, gateway/model evidence, and OS runtime telemetry as separate evidence channels; materializes them into one typed execution graph; and makes cross-layer attribution conservative, inspectable, and explicitly fallible. The design prefers an unattributed effect to a convenient but unsupported edge and separates evidence strength from behavioral severity.

This formulation positions \system between two mature traditions: application/distributed tracing, which excels at instrumented logical workflows, and whole-system provenance, which excels at machine-level dependency histories. Agentic software increasingly needs both. A submission-quality evaluation must therefore measure not only how many nodes a dashboard can display, but whether the system reconstructs the right entities and relations, abstains when evidence is ambiguous, remains consistent across providers and operating systems, survives adverse execution conditions, imposes acceptable cost, and materially improves human debugging and accountability tasks. The first manuscript establishes that model and the controlled evaluation protocol; the next empirical revision should fill those questions with frozen, reproducible measurements rather than product-feature inference.
